\chapter{Introduction}
\label{chap:intro}
\pagenumbering{arabic}

Supply chain management has become a major driver of business performance as logistic costs are on a continuous rise \citep{kearney2026geopolitical}. Therefore, it is imperative for firms to optimize their internal operations to remain competitive. One high impact area to do so is \gls{acr:ic} as inventory expenses represent a considerable share of a firms’ total logistics costs \citep{roodbergen2014simultaneous, lancioni2000new}. 

The central question in \gls{acr:ic} is how much each warehouse should reorder so that the cost of holding stock is balanced against the cost of stocking out. This is non-trivial since demand is stochastic, replenishment arrives after a lead time, and a single warehouse might manage many products simultaneously. The difficulty grows in real distribution networks, where warehouses replenish stock independently from local information and differ widely in the conditions they face. On balance, decentralization and heterogeneity are the central challenges in real-world \gls{acr:ic} and thus one of the largest cost-drivers. 

Replenishment has long been guided by classical models such as the newsvendor and base-stock policies, which compute order quantities based on assumptions about demand and costs \citep{scarf1960sSoptimality, Chen2026}. These models are tractable and well understood, but the assumptions behind their solution guarantees often fail in practice. As a response, \gls{acr:rl} has recently gained popularity since it learns policies directly from interaction and handles uncertainty that the classical models cannot \citep{boute2022roadmap}. Most of this work however assumes a single central controller that has access to the whole network and places collective orders, which seldom reflects how firms operate in practice.

To address this, \gls{acr:marl} has been proposed as the more realistic alternative. It models each warehouse as a separate agent that acts on local information and thus has to interact with other agents in the network. Promisingly, cooperative \gls{acr:marl} with \gls{acr:ps} already has been shown to perform well on multi-warehouse inventory problems \citep{mousa2024analysis}. Most \gls{acr:marl} approaches however have in common that they act on symmetric environment instances. Since real networks are not symmetric and have multiple warehouses with different demand profiles, non-uniform products, and asymmetric cost regimes, it is still unclear how parameter-shared \gls{acr:marl} behaves once heterogeneity is introduced. Consequently a single shared policy must then act differently for each warehouse while identifying them only from the local observation, which is exactly where it may fail.

Therefore, this thesis studies cooperative \gls{acr:marl} with \gls{acr:ps} for inventory control under heterogeneity in demand, product portfolio, and shipment
costs. We organize the study around two central research questions:

\paragraph{\textbf{RQ1}}
Under heterogeneity, does \gls{acr:ppo} for the multi-agent case with \gls{acr:ps} degrade relative to a symmetric baseline of matched network throughput? If so, what failure modes explain the resulting gap and can environment-specific hyperparameter tuning recover lost performance?

\vspace{-0.7\baselineskip}
\paragraph{\textbf{RQ2}}
\gls{acr:ps} imposes a single inductive bias on every agent that becomes a
liability under heterogeneity. Which architectural modifications relax this
constraint enough to recover any lost performance and is one modification effective
across all heterogeneity types or does each type require a distinct architectural
response?\\

To answer these questions, we build a configurable environment with controlled
heterogeneity and an order-level demand model for multiple \glspl{acr:sku} that lets us test our architectural interventions against the failure modes we identify. In doing so, we provide three main contributions:

\paragraph{\textbf{C1}} An extensible benchmark that brings \gls{acr:marl} inventory control closer to operational practice by contributing a simulation environment that supports controlled heterogeneity across multiple network dimensions and adopts an order level multi-\gls{acr:sku} demand model.

\vspace{-0.7\baselineskip}
\paragraph{\textbf{C2}} A diagnostic characterization of the performance degradation of cooperative \gls{acr:marl} with \gls{acr:ps} across heterogeneity types as well as the underlying failure modes to show where and how standard \gls{acr:marl} breaks in heterogeneous settings.

\vspace{-0.7\baselineskip}
\paragraph{\textbf{C3}} A toolkit of architectural interventions to address the diagnosed failure modes and an analysis of both their individual efficacy and the effects of combining them.\\

The remainder of the thesis is organized as follows. \autoref{chap:background-and-related-work} reviews related work and derives any prevalent gaps. \autoref{chap:problem-formulation} then formalizes the problem setting that we will use for our experiments. While \autoref{chap: background} proceeds to outline theoretical foundations in \gls{acr:rl}, \autoref{chap: methodology} applies \gls{acr:rl} methods to the distribution network of \autoref{chap:problem-formulation} and presents the observations, actions, and algorithms.  Finally, \autoref{chap:experiment-results} defines the experimental setup and reports the experimental results. The thesis concludes with a conclusion of the results in \autoref{chap:conclusion}.



